\section{Podstawy teoretyczne}
\label{sec:teoria}
\subsection{Kwalifikacja pomieszczenia}
\todo{Opisać wyznaczanie chłonności akustycznej z czasu pogłosu i objętości
pomieszczenia oraz metodą przybliżoną ze współczynników pochłaniania.}
\todo{Podać zależność poprawki środowiskowej od chłonności i powierzchni
pomiarowej oraz kryterium kwalifikacji. Rozróżnić wyniki pasmowe i A.}

\subsection{Wyznaczanie poziomu mocy akustycznej}
\todo{Opisać energetyczne uśrednianie poziomów na powierzchni, poprawkę
na tło, poprawkę środowiskową i przeliczenie poziomu powierzchniowego
na poziom mocy akustycznej. Zdefiniować symbole i wielkości odniesienia.}
Podstawę opracowania stanowią oba konspekty
\autocite{konspektKwalifikacja,konspektMoc}.
